

\cleardoublepage
\thispagestyle{empty}
\bigtitle{In Octava Sanct. Cord. Jesu.}
\addcontentsline{toc}{section}{XIX Junii} \phantomsection
\duplexmajus

\header{xix junii 2026}

\office{ad primam}
\addcontentsline{toc}{subsection}{Ad Primam} \phantomsection

\pateravecredo

\gscore[]{℣.}{Deus_in_adiutorium_Tonus_ferialis_no_laus_tibi}

\gscore[Hy.]{3.}{hy_jam_lucis_orto_sidere_in_festo_s._cordis_jesu_solesmes_1961}

\gscore[Ant.]{1 f.}{unus_militum}
\psalmus{53}{53_1f}{53_1}
\psalmus{118.\textsc{i}}{118_1_1f}{118_1_1}
\psalmus{118.\textsc{ii}}{118_2_1f}{118_2_1}
\smallscore{an_unus_militum}

\capitulum{1 Tim. 1, 17}
\lettrine{R}{e}gi sæculórum immortáli et invisíbili,~† soli Deo honor et glória~* in sǽcula
sæculórum.

\rr Deo grátias.

\gscore[]{Resp. brev.}{rb_christe_fili_dei_vivi_sc}

\smallscore{versiculus_ad_primam}

  \vv Dóminus vobíscum. \rr Et cum spíritu tuo.
  
  \oratio
  
  Orémus.
  
 \lettrine{D}{o}mine Deus omnípotens, qui ad princípium huius diéi nos perveníre
fecísti: tua nos hódie salva virtúte; ut in hac die ad nullum declinémus peccátum,
sed semper ad tuam justítiam faciéndam nostra procédant elóquia, dirigántur
cogitatiónes et ópera. Per Dóminum.

\rr Amen.

\vv Dóminus vobíscum. \rr Et cum spíritu tuo.

\smallscore{BD_ad_completorium}

\smalltitle{Martyrologium.}

\begin{enpars}
\lettrine[lraise=0.2]{J}{u}ne 20th 2026, the fifth day of the Moon, were born into the better life:
The holy martyr Pope Silverius. He refused to restore the (Eutychian) heretic Anthimus, who had been deposed (from the Patriarchate of Constantinople) by his predecessor, Pope Agapitus, and in consequence of this, through a plot of the wicked Empress Theodora, was exiled by Bellisarius to the island of Ponza, where he died for the Catholic faith, broken down by sufferings and hardships, (in the year 538.)
At Rome, the holy Novatus, son of the blessed Senator Pudens and brother of the holy Priest Timothy, and of Christ's holy Virgins Pudentiana and Praxedes, who were taught in the faith by the Apostles. Their house was turned into a church, and called that of St. Pastor, (in the second century.)
At Tomi, in Pontus, the holy martyrs Paul and Cyriacus.
At Petra, in Palestine, the holy Macarius, Bishop (of that see) who suffered many things of the Arians, and was exiled to Africa, where he fell asleep in the Lord, (in the fourth century.)
At Seville, in Spain, (in the year 630,) the holy Virgin Florentina, sister of the holy Leander, Bishop (of Seville,) and holy Isidore, Bishop (of Seville.)
\end{enpars}

\vv Et álibi aliórum plurimórum san\-ctórum Mártyrum et Confessórum, atque san\-ctárum Vírginum.

\rr Deo grátias.

\vv Pretiósa in conspéctu Dómini. \rr Mors Sanctórum ejus.

Sancta María et omnes Sancti intercédant pro nobis ad Dóminum, ut nos
mereámur ab eo adjuvári et salvári, qui vivit et regnat in sǽcula sæculórum.

\rr Amen.

\vv Deus in adjutórium meum inténde.

\rr Dómine, ad adjuvándum me festína.

\vv Deus in adjutórium meum inténde.

\rr Dómine, ad adjuvándum me festína.

\vv Deus in adjutórium meum inténde.

\rr Dómine, ad adjuvándum me festína.

\vv Glória Patri, et Fílio,~* et Spirítui Sancto.

\rr Sicut erat in princípio, et nunc, et semper,~* et in sǽcula sæculórum. Amen.

\vv Kýrie eléison.

\rr Christe eléison. Kýrie eléison.

\vv Pater noster \textit{secreto.}

\vv Et ne nos indúcas in tentatiónem:

\rr Sed líbera nos a malo.

\vv Réspice in servos tuos, Dómine, et in ópera tua, et dírige fílios eórum.

\rr Et sit splendor Dómini Dei nostri super nos, et ópera mánuum nostrárum
dírige super nos, et opus mánuum nostrárum dírige.

\vv Glória Patri, et Fílio,~* et Spirítui Sancto.

\rr Sicut erat in princípio, et nunc, et semper,~* et in sǽcula sæculórum. Amen.

Orémus.

\lettrine{D}{i}rígere et sanctificáre, régere et gubernáre dignáre, Dómine Deus, Rex
cæli et terræ, hódie corda et córpora nostra, sensus, sermónes et a\-ctus nostros in
lege tua, et in opéribus mandatórum tuórum: ut hic et in ætérnum, te auxiliánte,
salvi et líberi esse mereámur, Salvátor mundi: Qui vivis et regnas in sǽcula
sæculórum.

\rr Amen.

\vv Jube, domne, benedícere.

\smalltitle{Benedictio.}

Dies et actus nostros in sua pace dispónat Dóminus omnípotens.

\rr Amen.

\lectiobrevis{Eph. 3, 17-19}
\lettrine{I}{n} caritáte radicáti et fundáti, ut possítis comprehéndere cum ómnibus sanctis, quæ sit latitúdo et longitúdo, et sublímitas et profúndum; scire étiam supereminéntem sciéntiæ caritátem Christi, ut impleámini in omnem plenitúdinem Dei.

\vv Tu autem, Dómine, miserére nobis.

\rr Deo grátias.

\vv Adjutórium nostrum in nómine Dómini.

\rr Qui fecit cælum et terram.

\vv Benedícite.

\rr Deus.

\lettrine{D}{o}minus nos benedícat, \cc{} et ab omni malo deféndat, et ad vitam perdúcat ætérnam. Et fidélium ánimæ per misericórdiam Dei requiéscant in pace.

\rr Amen.

\secreto
\office{ad missam}
\addcontentsline{toc}{subsection}{Ad Missam} \phantomsection
\proper{Introitus}{Ps. 32, 11, 19; 1}

\texteparacol{\lettrine{C}{o}gitatiónes Cordis ejus in generatióne et generatiónem: ut éruat a morte ánimas eórum et alat eos in fame.
℣. Exsultáte, justi, in Dómino: rectos decet collaudátio.
Glória Patri. Sicut erat.
Cogitatiónes.}{english}{\noindent The thoughts of his heart are to all generations: to deliver them from death and preserve them in spite of famine. ℣. Exult, you just, in the Lord; praise from the upright is fitting. Glory be to the Father. As it was in the beginning. The thoughts.}

\masspart{Oratio}

\texteparacol{\lettrine{D}{e}us, qui nobis in Corde Fílii tui, nostris vulneráto peccátis, infinítos dile\-ctiónis thesáuros misericórditer largíri dignáris:~† concéde quǽsumus; ut illi devótum pietátis nostræ præstántes obséquium,~* dignæ quoque satisfa\-ctiónis exhibeámus offícium. Per eúmdem Dóminum.

\lettrine[ante=2]{D}{e}us, qui beátam Juliánam Vírginem tuam extrémo morbo laborántem, pretióso Fílii tui Córpore mirabíliter recreáre dignátus es:~† concéde, quǽsumus; ut ejus intercedéntibus méritis, nos quoque eódem in mortis agóne refé\-cti ac roboráti,~* ad cæléstem pátriam perducámur.

\lettrine[ante=3]{D}{e}us, qui nos ánnua san\-ctórum Mártyrum túorum Gervásii et Protásii sole\-mnitáte lætíficas:~† concéde propítius; ut quorum gaudémus méritis,~* accendámur exémplis. Per Dóminum.}{english}{\noindent O God, who hast suffered the heart of thy Son to be wounded by our sins, and in that very heart hast bestowed on us the abundant riches of thy love: grant that the devout homage of our hearts, which we render unto him, may of thy mercy be deemed a recompence acceptable in thy sight. Through the same Lord.

\noindent 2. O God, who, when thy blessed handmaiden Juliana was lying sick unto death wast pleased in wondrous wise to comfort her with the Precious Body of thy Son, be thou entreated for the same thy servant's sake, and grant unto us also the same comfort in our last agony, that we may go in the strength of that Meat unto our very fatherland, which is in heaven.

\noindent 3. O God, who gladden us by the annual festival of thy martyrs Gervasii and Protasi, grant that we may be inspired by the example of those in whose merits we rejoice.
Through our Lord.}

\reading{Epistola}{Eph. 3, 8 -- 12, 14 -- 19}

\texteparacol{\lettrine{F}{r}atres: Mihi, ómnium san\-ctórum mínimo, data est grátia hæc, in géntibus evangelizáre investigábiles divítias Christi, et illumináre o\-mnes, quæ sit dispensátio sacraménti abscónditi a sǽculis in Deo, qui ó\-mnia creávit: ut innotéscat principátibus et potestátibus in cæléstibus per Ecclésiam multifórmis sapiéntia Dei, secúndum præfinitiónem sæculórum, quam fecit in Christo Jesu, Dómino nostro, in quo habémus fidúciam et accéssum in confidéntia per fidem ejus. Hujus rei grátia fle\-cto génua mea ad Patrem Dómini nostri Jesu Christi, ex quo omnis patérnitas in cælis ei in terra nominátur, ut det vobis, secúndum divítias glóriæ suæ, virtúte corroborári per Spíritum ejus in interiórem hóminem, Christum habitáre per fidem in córdibus vestris: in caritáte radicáti et fundáti, ut póssitis comprehéndere cum ó\-mnibus san\--ctis, quæ sit latitúdo, et longitúdo, et sublímitas, et profúndum: scire étiam supereminéntem sciéntiæ caritátem Christi, ut impleámini in omnem plenitúdinem Dei.}{english}{\noindent Brethren: To me, the very least of all saints, there was given this grace, to announce among the Gentiles the good tidings of the unfathomable riches of Christ, and to enlighten all men as to what is the dispensation of the mystery which has been hidden from eternity in God, Who created all things; in order that through the Church there be made known to the Principalities and the Powers in the heavens the manifold wisdom of God according to the eternal purpose which He accomplished in Christ Jesus our Lord. In Him we have assurance and confident access through faith in Him. For this reason I bend my knees to the Father of our Lord Jesus Christ, from Whom all fatherhood in heaven and on earth receives its name, that He may grant you from His glorious riches to be strengthened with power through His Spirit unto the progress of the inner man; and to have Christ dwelling through faith in your hearts: so that, being rooted and grounded in love, you may be able to comprehend with all the saints what is the breadth and length and height and depth, and to know Christ’s love which surpasses knowledge, in order that you may be filled unto all the fullness of God.}

\proper{Graduale \& Alleluia}{Ps. 24, 8 -- 9; Matt. 11, 29}

\texteparacol{\lettrine{D}{u}lcis et rectus Dóminus: propter hoc legem dabit delinquéntibus in via.
℣. Díriget mansúetos in judício, docébit mites vias suas.

\lettrine{A}{l}elúia, allelúia. Tóllite jugum meum super vos, et díscite a me, quia mitis sum et húmilis Corde, et inveniétis réquiem animábus vestris. Allelúia.}{english}{\noindent Good and upright is the Lord; thus He shows sinners the way.
℣. He guides the humble to justice; he teaches the humble his way.

\noindent Alleluia, alleluia. Take my yoke upon you, and learn from me, for I am meek, and humble of heart: and you will find rest for your souls. Alleluia.}

\reading{Evangelium}{Joann. 19, 31 -- 37}

\texteparacol{\lettrine{I}{n} illo témpore: Judǽi, quóniam Parascéve erat, ut non remanérent in cruce córpora sábbato, erat enim magnus dies ille sábbati, rogavérunt Pilátum, ut frangeréntur eórum crura, et tolleréntur. Venérunt ergo mílites: et primi quidem fregérunt crura et alteríus, qui crucifíxus est cum eo. Ad Jesum autem cum veníssent, ut vidérunt eum jam mórtuum, non fregérunt ejus crura, sed unus mílitum láncea latus ejus apéruit, et contínuo exívit sanguis et aqua. Et qui vidit, testimónium perhíbuit: et verum est testimónium ejus. Et ille scit quia vera dicit, ut et vos credátis. Fa\-cta sunt enim hæc ut Scriptúra implerétur: Os non comminuétis ex eo. Et íterum alia Scri\-ptúra dicit: Vidébunt in quem transfixérunt.}{english}{\noindent At that time, the Jews, since it was the Preparation Day, in order that the bodies might not remain upon the cross on the Sabbath,  for that Sabbath was a solemn day, besought Pilate that their legs might be broken, and that they might be taken away. The soldiers therefore came and broke the legs of the first, and of the other, who had been crucified with him. But when they came to Jesus, and saw that he was already dead, they did not break his legs; but one of the soldiers opened his side with a lance, and immediately there came out blood and water. And he who saw it has borne witness, and his witness is true; and he knows that he tells the truth, that you also may believe. For these things came to pass that the Scripture might be fulfilled: ``Not a bone of him shall you break.'' And again another Scripture says: ``They shall look upon him whom they have pierced.''}

\proper{Offertorium}{Ps. 68, 21}

\texteparacol{\lettrine{I}{m}propérium exs\-pectávi Cor meum et misériam: et sustínui, qui simul mecum contristarétur, et non fuit: consolántem me quæsívi, et non invéni.}{english}{\noindent My heart expected reproach and misery; I looked for sympathy, but there was none; and for comforters, and I found none.}

\masspart{Secreta}
\texteparacol{\lettrine{R}{e}spice, quǽsumus, Dómine, ad ineffábilem Cordis dilé\-cti Fílii tui caritátem: ut quod offérimus sit tibi munus accé\-ptum et nostrórum expiátio deli\-ctórum.
Per eúmdem Dóminum.

\lettrine[ante=2]{A}{c}cépta tibi sit, Dómine, sacrátæ plebis oblátio pro tuórum honóre San\-ctórum: quorum se méritis de tribulatióne percepísse cognóscit auxílium.

\lettrine[ante=3]{O}{b}látis, quǽsumus, Dómine, placáre munéribus: et, intercedéntibus san\-ctis Martýribus tuis Gervásio et Protásio, a cun\-ctis nos defénde perículis.
Per Dóminum.}{english}{\noindent Look, we beseech thee, o Lord, upon the heart of thy beloved Son, with its boundless love, so that what we offer, may be an acceptable gift in thy sight and an atonement for our sins.
Through the same Lord.

\noindent 2. May the service of thy holy people be acceptable to thee, o Lord, and honorable to thy saints, through whose merits we know that we have received help in trouble.

\noindent 3. Be appeased by the gifts we offer thee, o Lord, and through the intercession of thy holy martyrs Gervasii and Protasi, safeguard us from all dangers.
Through our Lord.}

\masspart{Præfatio de sacratissimo Cordis Jesu}

\texteparacol{\lettrine{V}{e}re dignum et justum est, æquum et salutáre, nos tibi semper et ubíque grátias ágere: Dómine san\-cte, Pater omnípotens, ætérne Deus: Qui Unigénitum tuum, in Cruce pendéntem, láncea mílitis transfígi voluísti: ut apértum Cor, divínæ largitátis sacrárium, torréntes nobis fúnderet miseratiónis et grátiæ: et, quod amóre nostri flagráre numquam déstitit, piis esset réquies et pæniténtibus pater et salútis refúgium. Et ídeo cum Angelis et Archángelis, cum Thronis et Dominatiónibus cumque omni milítia cæléstis exércitus hymnum glóriæ tuæ cánimus, sine fine dicéntes:}{english}{\noindent It is truly meet and just, right and for our salvation, that we should at all times, and in all places, give thanks unto Thee, O holy Lord, Father almighty, everlasting God; Who didst will that Thine only-begotten Son, while hanging on the cross, should be pierced by a soldier's spear, that the Heart thus opened, a shrine of divine bounty, should pour out on us streams of mercy and grace, and that what never ceased to burn with love for us, should be a resting-place to the devout, and open as a refuge of salvation to the penitent. And therefore with angels and archangels, with thrones and dominions, and with all the hosts of the heavenly army, we sing the hymn of thy glory, evermore saying:}

\proper{Communio}{Joann. 19, 34}

\texteparacol{\lettrine{U}{n}us mílitum láncea latus ejus apéruit, et contínuo exívit sanguis et aqua.}{english}{\noindent One of the soldiers opened his side with a lance, and immediately there came out blood and water.}

\masspart{Postcommunio}
\texteparacol{\lettrine{P}{r}ǽbeant nobis, Dómine Jesu, divínum tua san\-cta fervórem: quo dulcíssimi Cordis tui suavitáte percé\-pta; discámus terréna despícere, et amáre cæléstia:
Qui vivis et regnas.

\lettrine[ante=2]{S}{a}tiásti, Dómine, famíliam tuam munéribus sacris: ejus, quǽsumus, semper interventióne nos réfove, cujus sollémnia celebrámus.

\lettrine[ante=3]{H}{æ}c nos commúnio, Dómine, purget a crímine: et, intercedéntibus san\-ctis Martýribus tuis Gervásio et Protásio, cæléstis remédii fáciat esse consórtes. Per Dóminum.}{english}{\noindent May thy sacrament, o Lord Jesus, give us holy zeal, so that, seeing the sweetness of thy most loving heart, we may learn to despise the things of earth and love those of heaven.
Who livest and reignest.

\noindent 2. Thou hast fed thy servants, o Lord, with the holy gifts; comfort us ever we pray, by her intercession whose festival we celebrate.

\noindent 3. May this communion, o Lord, cleanse us of sin, and, through the intercession of thy holy martyrs Gervasii and Protasi, make us sharers of heavenly salvation.
Through our Lord.}

\office{ad vesperas}
\addcontentsline{toc}{subsection}{Ad  Vesperas} \phantomsection
%\label{in_2_vesperis_sc}\phantomsection

\gscore[]{℣.}{Deus_In_Adjutorium_Festivus}

\gscore[1. Ant.]{1 f.}{an_unus_militum_solesmes_1961}
\psalmus{109}{109_1f}{109_1}
\smallscore{an_unus_militum}

\gscore[2. Ant.]{7 c.}{an_stans_jesus_solesmes_1961} %% LaTeX's tracking is HORRIBLE with this psalm. Lots of white for no reason. %%needs fixing
\psalmus{110}{110_7c}{110_7}
\smallscore{an_stans_jesus}

\gscore[3. Ant.]{3. a2}{an_in_caritate_solesmes_1961}
\psalmus{115}{115_3a2}{115_3a2_g}
\smallscore{an_in_caritate}

\gscore[4. Ant.]{4. E}{an_venite_ad_me_solesmes_1961}
\psalmus{127}{127_4E}{127_4E}
\smallscore{an_venite_ad_me}

\gscore[5. Ant.]{5. a}{an_fili_praebe_solesmes_1961}
\psalmus{147}{147_5}{147_5}
\smallscore{an_fili_praebe}
\capitulum{Eph. 3, 8 – 9}

\lettrine{F}{r}ratres: Mihi ó\-mnium san\-ctórum mínimo data est grátia hæc,~† in géntibus evangelizáre investigábiles divítias Christi:~* et illumináre omnes, quæ sit dispensátio sacraménti abscónditi a sǽculis in Deo.

\rr Deo grátias.

\hymnus
   \gscore[]{3.}{hy_en_ut_superba_solesmes_1961}
   
   \smallscore{versiculus_haurietis_tonus_solius}

   \admagnificat
   
   \gscore[]{1.f}{an_ad_jesum_autem_solesmes_1961}
    
    \smallscore{Magnificat_1f}
    
  \magnificattexte{Magnificat_1_ordinaire}
  
  \smallscore{an_ad_jesum_autem}
  
  \vv Dóminus vobíscum. \rr Et cum spíritu tuo.
  
  \oratio
  
  Orémus.
  
  \lettrine{D}{e}us, qui nobis in Corde Fílii tui, nostris vulneráto peccátis, infinítos dile\-ctiónis thesáuros misericórditer largíri dignáris:~† concéde quǽsumus; ut illi devótum pietátis nostræ præstántes obséquium,~* dignæ quoque satisfa\-ctiónis exhibeámus offícium. Per eúmdem Dóminum.
  
  \rr Amen.
  
 \gscore[Ant.]{2.}{an_beata_mater_in_sabbato_solesmes_1961}
 
 \vv Diffúsa est grátia in lábiis tuis.

\rr Propteréa benedíxit te Deus in ætérnum.

  Orémus.
  
  \lettrine{C}{o}ncéde nos famúlos tuos, quæsumus Domíne Deus, perpétua mentis et corpóris sanitáte gaudére:~† et gloriósa beátæ Mariæ semper Virgínis intercessióne,~* a præsénti liberári tristítia, et æterna perfrúi lætítia.
  
\gscore[Ant.]{7.}{an_cs_veni_sponsa_ii_vesp_solesmes_1961}

\vv Spécie tua et pulchritúdine tua.

\rr Inténde, próspere procéde, et regna.

Orémus.

\lettrine{D}{e}us, qui beátam Juliánam Vírginem tuam extrémo morbo laborántem, pretióso Fílii tui Córpore mirabíliter recreáre dignátus es:~† concéde, quǽsumus; ut ejus intercedéntibus méritis, nos quoque eódem in mortis agóne refé\-cti ac roboráti,~* ad cæléstem pátriam perducámur.

\gscore[Ant.]{8.}{an_cs_iste_sanctus_pro_lege_solesmes_1961}

\vv Glória et honóre coronásti eum Dómine.

\rr Et constituísti eum super ópera mánuum tuárum.

\newpage

Orémus.

%\lettrine{I}{n}firmitátem nostram réspice, omnípotens Deus:~† et quia pondus própriæ actiónis gravat,~* beáti Silvérii Mártyris tui atque Pontíficis intercéssio gloriósa nos prótegat. Per Dóminum.

\lettrine{G}{r}egem tuum, Pastor ætérne, placátus inténde:~† et per beátum Silvérium Mártyrem tuum atque Summum Pontíficem, perpétua prote\-ctióne custódi;~* quem totíus Ecclésiæ præstitísti esse pastórem. Per Dóminum.

\rr Amen.

   \vv Dóminus vobíscum. \rr Et cum spíritu tuo.
   
\gscore[]{8.}{BD_Duplex_II_Vesp}
   
  \vv Fidélium ánimæ per misericórdiam Dei requiéscant in pace.
  
  \rr Amen.
  
 \rubrique{Tunc dicitur Completorium. Si non dicendum:}
  
  \secreto
  
\vv Dóminus det nobis suam pacem.
  
\rr Et vitam ætérnam. Amen.

\rubrique{\normaltext{Salve Regína,} etc., \normaltext{\pageref{an_salve_regina_solesmes}.}}